% Authored: architecture decision records.

\chapter{Architecture Decision Records}
\label{ch:adr}

\section*{Method}
\addcontentsline{toc}{section}{Method}

The records below document decisions that are \emph{observable} in the
repository at revision \texttt{\projectrevision}: a structural choice that is
consistently applied and that could plausibly have been made differently.

Two rules were applied while writing them.

\begin{enumerate}
  \item A decision is recorded only when the repository contains evidence for
    it: the code itself, a source comment or docstring that states the reason,
    a configuration file, a CI step, or the contributor instructions in
    \srcfile{AGENTS.md}.
  \item Where the \emph{rationale} cannot be established from repository
    evidence, the record says so explicitly instead of supplying a plausible
    motivation. The repository has no ADR directory, no design document, and no
    commit messages that argue for a design; \srcfile{CHANGELOG.md} records
    what changed, not why.
\end{enumerate}

Every record therefore carries an \textbf{Evidence} subsection naming the files
that support it. Status values are descriptive of the current revision: an
\emph{Accepted} record means the decision is in force in the source, not that a
review body approved it.

\section{ADR-001: Every default is local and deterministic}
\label{adr:001}

\subsection*{Status}

Accepted; in force at revision \texttt{\projectrevisionshort}.

\subsection*{Context}

The system's purpose requires a language model, an embedding model, and a
vector store. Each of those is normally a remote or heavyweight dependency,
which would make the repository impossible to test, demonstrate, or run in CI
without credentials and network access.

\subsection*{Decision}

Every component that would normally require external infrastructure has a
local deterministic default, and the default configuration selects it:

\begin{itemize}[nosep]
  \item \texttt{FEEDBACK\_AGENT\_LLM\_PROVIDER=local} selects
    \texttt{DeterministicLLM};
  \item embeddings default to \texttt{HashingEmbeddingModel};
  \item \texttt{FEEDBACK\_AGENT\_VECTOR\_STORE=json} selects
    \texttt{InMemoryVectorStore} with JSON persistence;
  \item reranking, query rewriting, and hallucination checking all default to
    deterministic local implementations;
  \item telemetry defaults to a no-op emitter.
\end{itemize}

\subsection*{Consequences}

\begin{itemize}
  \item The full pipeline --- ingest, chunk, embed, index, retrieve, rerank,
    generate, parse, cite, check --- runs offline. CI executes
    \srcfile{scripts/run\_demo.py} on every push as a blocking step, so
    breaking the offline path breaks the build.
  \item Test assertions can be exact rather than fuzzy, because generation is
    deterministic. The prompt snapshot test compares rendered output byte for
    byte.
  \item Evaluation metrics are reproducible across runs, which is what makes
    them usable as a regression gate.
  \item The deterministic provider is not a language model. Answer quality
    under the default configuration reflects the retrieval and templating
    pipeline only, and evaluation numbers obtained with it do not transfer to a
    remote provider.
  \item A second code path exists for every capability, and both must be kept
    working.
\end{itemize}

\subsection*{Alternatives}

Requiring an API key for any use; recording and replaying provider fixtures;
mocking at the test boundary only. The first two would leave CI unable to
exercise the real pipeline; the third would leave the demo and the CLI
unusable without credentials.

\subsection*{Evidence}

\srcfile{src/feedback\_intelligence\_agent/config.py} (default field values);
\srcfile{src/feedback\_intelligence\_agent/factory.py} (\texttt{build\_llm},
\texttt{build\_retriever}, \texttt{load\_or\_build\_index});
\srcfile{.github/workflows/ci.yml} (``Run deterministic demo'' step);
\srcfile{AGENTS.md} (``Keep the local deterministic path working without
external APIs'' and ``The demo must continue to work without an API key'');
\srcfile{tests/test\_prompt\_snapshot.py}.

\section{ADR-002: Substitution through structural protocols, not inheritance}
\label{adr:002}

\subsection*{Status}

Accepted; in force at revision \texttt{\projectrevisionshort}.

\subsection*{Context}

Every swappable capability --- retrieval, generation, embedding, reranking,
query rewriting, answer judging, vector storage, stream consumption --- needs a
stable interface, and the local and remote implementations of each have nothing
useful to share by inheritance.

\subsection*{Decision}

Interfaces are declared as \texttt{typing.Protocol} classes and implementations
do not inherit from them. The protocols are \texttt{Retriever},
\texttt{LLMProvider}, \texttt{EmbeddingModel}, \texttt{Reranker},
\texttt{QueryRewriter}, \texttt{HallucinationJudge}, \texttt{VectorStore}, and
\texttt{FeedbackStream}. Only \texttt{tools.Tool} uses an abstract base class,
because tools share concrete payload-building behaviour.

\subsection*{Consequences}

\begin{itemize}
  \item A test double satisfies an interface by having the right methods; the
    suite uses plain local classes rather than mocking frameworks.
  \item \texttt{mypy} in strict mode checks conformance structurally, so an
    implementation that drifts from the protocol fails the type gate rather
    than failing at run time.
  \item Optional SDK objects can satisfy an interface without the package
    importing the SDK at module level.
  \item Protocols carry no shared implementation, so genuinely common
    behaviour --- the resilience policy, for example --- has to be added by
    composition (\texttt{ResilientLLMProvider} wraps a provider) rather than
    inherited.
  \item \texttt{VectorStore} is additionally \texttt{@runtime\_checkable}, so
    it supports \texttt{isinstance} checks, which the other protocols do not.
\end{itemize}

\subsection*{Alternatives}

Abstract base classes for every interface, which would force optional adapters
to import the package to subclass it, and would make third-party objects
non-conforming even when they have the right shape.

\subsection*{Evidence}

\srcfile{src/feedback\_intelligence\_agent/retrieval.py},
\srcfile{embeddings.py}, \srcfile{llm.py}, \srcfile{reranking.py},
\srcfile{memory.py}, \srcfile{hallucination.py}, \srcfile{vector\_store.py},
\srcfile{streaming\_ingestion.py} (protocol definitions);
\srcfile{src/feedback\_intelligence\_agent/tools.py} (the one ABC);
\srcfile{pyproject.toml} (\texttt{[tool.mypy] strict = true});
\srcfile{AGENTS.md} (``Prefer explicit interfaces over hidden global state'').

\section{ADR-003: One composition root}
\label{adr:003}

\subsection*{Status}

Accepted; in force at revision \texttt{\projectrevisionshort}.

\subsection*{Context}

Three entry points --- library, CLI, and HTTP API --- need the same wired
object graph, and that graph depends on configuration in several independent
dimensions: provider, retriever, vector store, telemetry backend, resilience,
and judge.

\subsection*{Decision}

All construction lives in
\cref{module:feedback_intelligence_agent.factory}. Every consumer calls a
\texttt{build\_*} function; no other module branches on a configuration value
to choose an implementation.

\subsection*{Consequences}

\begin{itemize}
  \item \texttt{factory} has the highest fan-out in the package, which the
    coupling table in \cref{ch:dependencies} shows directly. That is the
    intended shape of a composition root rather than an accident.
  \item Every other module can be constructed directly with explicit
    collaborators, which is what the unit tests do; they never go through the
    factory.
  \item Adding a provider is a two-file change: an adapter class and one branch
    in \texttt{build\_llm}.
  \item A configuration error surfaces at construction with a message naming
    the missing variable, not at first use.
  \item Because the factory imports nearly everything, importing it eagerly
    would pull in optional adapters; this is the direct cause of the deferred
    imports recorded in \cref{adr:005} and \cref{adr:006}.
\end{itemize}

\subsection*{Alternatives}

A dependency-injection container, which \srcfile{AGENTS.md} discourages (``Do
not introduce heavyweight frameworks''); or per-entry-point construction, which
would duplicate the provider-selection logic three times.

\subsection*{Evidence}

\srcfile{src/feedback\_intelligence\_agent/factory.py};
\srcfile{api.py} (\texttt{create\_app} calls \texttt{build\_*} only);
\srcfile{cli.py} (imports \texttt{build\_agent}, \texttt{build\_index},
\texttt{build\_retriever}, \texttt{build\_telemetry}, \ldots);
\srcfile{docs/metadata/module-dependencies.json} (fan-out).

\section{ADR-004: Validated types at every boundary}
\label{adr:004}

\subsection*{Status}

Accepted; in force at revision \texttt{\projectrevisionshort}.

\subsection*{Context}

Data enters from CSV files, HTTP request bodies, JSONL streams, JSON files on
disk, and language-model responses. All of these are untrusted and
loosely-shaped.

\subsection*{Decision}

Every value crossing a boundary is a Pydantic model, and constraints are
declared on the fields rather than checked imperatively: rating bounds,
minimum text length, embedding-dimension range, metric ranges in $[0,1]$,
non-negative counters, and cross-field rules such as
\texttt{min\_rating} $\leq$ \texttt{max\_rating}. Persisted state --- stored
conversations, jobs, reports, human feedback, active-learning state --- is
re-validated on load rather than trusted.

\subsection*{Consequences}

\begin{itemize}
  \item Invalid input fails at the boundary with a message naming the field, so
    interior code can assume well-formed data and needs no defensive checks.
  \item FastAPI derives request validation and the OpenAPI schema from the same
    models, so the HTTP contract and the domain contract cannot diverge.
  \item \texttt{schemas} acquires the highest fan-in in the package (26 of 43
    modules import it), which the coupling table shows. The shared vocabulary
    is a deliberate coupling.
  \item Pydantic reaches the foundation layer and therefore, transitively,
    every layer. Replacing it would be a repository-wide change.
  \item Model construction costs more than a dictionary. The repository accepts
    this; the benchmark harness measures the pipeline including that cost.
\end{itemize}

\subsection*{Alternatives}

Dataclasses with hand-written validation, or dictionaries with runtime
assertions. Both would require duplicating the HTTP schema separately from the
domain types.

\subsection*{Evidence}

\srcfile{src/feedback\_intelligence\_agent/schemas.py},
\srcfile{data\_contracts.py}, \srcfile{config.py}, \srcfile{evaluation.py},
\srcfile{memory.py}, \srcfile{jobs.py}, \srcfile{reports.py};
\srcfile{AGENTS.md} (``Use Pydantic models for external request and response
schemas''; ``Raise specific exceptions for data quality issues'');
\srcfile{pyproject.toml} (\texttt{plugins = ["pydantic.mypy"]}).

\section{ADR-005: Optional integrations behind extras and deferred imports}
\label{adr:005}

\subsection*{Status}

Accepted; in force at revision \texttt{\projectrevisionshort}.

\subsection*{Context}

Six integrations --- Anthropic, AWS Bedrock, Kafka/Kinesis, MLflow,
OpenTelemetry, Qdrant --- bring large dependency trees. Requiring them would
make the default install heavy and would couple the package's supported Python
range to theirs.

\subsection*{Decision}

Each is declared as an optional Poetry dependency exposed through a named
extra (\texttt{anthropic}, \texttt{bedrock}, \texttt{streaming},
\texttt{mlflow}, \texttt{otel}, \texttt{qdrant}), imported inside the function
that needs it rather than at module level, and selected by a configuration
value. \texttt{mypy} carries a per-module \texttt{ignore\_missing\_imports}
override for each, so the type gate passes whether or not the extra is
installed.

\subsection*{Consequences}

\begin{itemize}
  \item \texttt{poetry install} without extras produces a working system, and
    CI does exactly that.
  \item A missing extra surfaces as an \texttt{ImportError} at the moment the
    feature is first used, not at import time --- later than a module-level
    import would fail, and only on the path that needs it.
  \item The optional code paths are not exercised against the real SDKs in CI.
    The suite covers them through injected fakes, so an upstream SDK change can
    break an adapter without any gate noticing.
  \item Each override in \srcfile{pyproject.toml} carries a comment explaining
    why it exists, which keeps the exception list auditable.
\end{itemize}

\subsection*{Alternatives}

Hard dependencies on everything; or separate distribution packages per
integration. The first inflates the install; the second is disproportionate for
six thin adapters.

\subsection*{Evidence}

\srcfile{pyproject.toml} (\texttt{[tool.poetry.extras]} and the six commented
\texttt{[[tool.mypy.overrides]]} blocks);
\srcfile{src/feedback\_intelligence\_agent/llm.py},
\srcfile{qdrant\_store.py}, \srcfile{streaming\_ingestion.py},
\srcfile{telemetry.py}, \srcfile{experiments.py} (function-local imports);
\srcfile{.github/workflows/ci.yml} (\texttt{poetry install -{}-with dev}, no
extras).

\section{ADR-006: Break the factory/jobs cycle with a deferred import}
\label{adr:006}

\subsection*{Status}

Accepted; in force at revision \texttt{\projectrevisionshort}.

\subsection*{Context}

\pkg{jobs} needs \texttt{chunk\_to\_embedding\_text} from \pkg{factory} so that
an ingestion job produces vectors identical to a full rebuild. \pkg{factory}
needs \texttt{JsonJobStore} from \pkg{jobs} so that \texttt{build\_job\_store}
can return one. At module level this is a runtime import cycle.

\subsection*{Decision}

\texttt{factory.build\_job\_store} imports \texttt{jobs} inside the function
body, and the annotation-only reference is satisfied by a
\texttt{TYPE\_CHECKING} import at module level. The same pattern is applied to
\texttt{build\_report\_store}, \texttt{build\_human\_feedback\_store}, and
\texttt{build\_active\_learning\_state\_store}.

The reason is stated in the source: \texttt{build\_job\_store} carries the
docstring line ``Imported lazily to avoid a circular import, since
\texttt{feedback\_intelligence\_agent.jobs} reuses
\texttt{chunk\_to\_embedding\_text}.'' This is the only decision in this
chapter for which the repository states its own rationale in prose.

\subsection*{Consequences}

\begin{itemize}
  \item The eager import graph is acyclic, which the generated cycle analysis
    confirms: zero strongly connected components among module-level imports.
  \item Counting deferred edges, exactly one component becomes cyclic:
    \texttt{\{factory, jobs\}}. \cref{ch:dependencies} reports it rather than
    hiding it, because the acyclicity is a property of import placement rather
    than of the design.
  \item The annotation remains type-checked despite the runtime import being
    absent.
  \item The cycle persists as a latent constraint: moving that import to module
    level at any future point reintroduces it.
\end{itemize}

\subsection*{Alternatives}

Extracting \texttt{chunk\_to\_embedding\_text} into a shared module that both
import, which would eliminate the cycle rather than defer it. The repository
records no consideration of this.

\noevidence

\subsection*{Evidence}

\srcfile{src/feedback\_intelligence\_agent/factory.py}
(\texttt{if TYPE\_CHECKING} block and the four function-local imports;
docstring of \texttt{build\_job\_store});
\srcfile{docs/metadata/module-dependencies.json}
(\texttt{cycles.runtime} is empty; \texttt{cycles.including\_deferred}
contains \texttt{["factory", "jobs"]}).

\section{ADR-007: File-backed JSON stores instead of a database}
\label{adr:007}

\subsection*{Status}

Accepted; in force at revision \texttt{\projectrevisionshort}.

\subsection*{Context}

Five kinds of state outlive a request: conversations, ingestion jobs, saved
insight reports, human answer feedback, and active-learning workflow state.

\subsection*{Decision}

Each is persisted as JSON files under a configurable directory beneath
\srcfile{.artifacts/}, one file per entity, through a small store class
(\texttt{JsonConversationStore}, \texttt{JsonJobStore},
\texttt{JsonInsightReportStore}, \texttt{JsonHumanFeedbackStore},
\texttt{JsonActiveLearningStateStore}). Records are re-validated against their
Pydantic model on load. No database, ORM, or migration tool is present in the
dependency set.

\subsection*{Consequences}

\begin{itemize}
  \item No service is required to run the API, and CI needs no database
    container.
  \item Stored state is inspectable and diffable with ordinary tools, which
    makes the workflow features easy to test.
  \item Each store is behind a class boundary, so replacing the backing store
    is a contained change.
  \item There is no transactionality and no cross-process locking. Concurrent
    writers to the same entity race on whole-file writes. The production
    Compose file runs multiple Gunicorn workers, which makes this reachable in
    that configuration.
  \item Listing is a directory scan, so it degrades linearly with the number of
    stored entities.
\end{itemize}

\subsection*{Alternatives}

SQLite, which would add transactionality at the cost of a schema and
migrations; or an external database, which would contradict \cref{adr:001}.
The repository records no evaluation of either.

\noevidence

\subsection*{Evidence}

\srcfile{src/feedback\_intelligence\_agent/memory.py}, \srcfile{jobs.py},
\srcfile{reports.py}, \srcfile{human\_feedback.py},
\srcfile{active\_learning.py};
\srcfile{src/feedback\_intelligence\_agent/config.py} (the five
\texttt{*\_store\_path} settings); \srcfile{pyproject.toml} (no database
dependency); \srcfile{deploy/docker-compose.prod.yml} (multiple workers).

\section{ADR-008: Two independent guardrail gates}
\label{adr:008}

\subsection*{Status}

Accepted; in force at revision \texttt{\projectrevisionshort}.

\subsection*{Context}

Two different untrusted texts reach the model in a retrieval-augmented flow:
the user's question, and the retrieved feedback, which is itself user-authored
content stored earlier. A single check on the question cannot protect against
instructions embedded in retrieved evidence.

\subsection*{Decision}

The guardrail is split into two gates with distinct entry points and distinct
placement in the flow:

\begin{enumerate}[nosep]
  \item \texttt{check\_input(question)} runs \emph{before} retrieval. A block
    short-circuits the whole request into a safe refusal --- no retrieval
    happens, no tool runs, no prompt is built, no provider is called.
  \item \texttt{check\_context(chunk\_texts)} runs \emph{after} retrieval and
    \emph{before} prompt construction. Suspicious chunks are dropped
    individually and the request continues with the survivors.
\end{enumerate}

\subsection*{Consequences}

\begin{itemize}
  \item A refused question costs no retrieval and no generation.
  \item Injection-shaped evidence degrades an answer rather than failing it:
    dropping chunks preserves availability.
  \item Both outcomes are observable. The refusal carries the
    \texttt{GuardrailDecision} on the answer with its reason and severity, and
    the number of dropped chunks appears in the run span and in the answer
    diagnostics under \texttt{guardrail\_context\_dropped}.
  \item The checks are deterministic pattern matching, so they run with no
    model call, but they are correspondingly limited: an injection phrased
    outside the pattern set passes.
\end{itemize}

\subsection*{Alternatives}

One combined check over question and context, which could not distinguish
``refuse the request'' from ``drop this evidence''; or a model-based classifier,
which would contradict \cref{adr:001}.

\subsection*{Evidence}

\srcfile{src/feedback\_intelligence\_agent/guardrails.py}
(\texttt{check\_input}, \texttt{check\_context},
\texttt{is\_suspicious\_context}, \texttt{SAFE\_REFUSAL});
\srcfile{src/feedback\_intelligence\_agent/agent.py}
(\texttt{answer} calls the input gate before \texttt{\_retrieve} and the
context gate after it); \srcfile{tests/test\_guardrails.py}.

\section{ADR-009: Stream by replaying a completed answer}
\label{adr:009}

\subsection*{Status}

Accepted; in force at revision \texttt{\projectrevisionshort}.

\subsection*{Context}

A streaming endpoint is expected of a chat-style API, but the default provider
is the deterministic local one, which has no token stream, and the pipeline
performs retrieval, reranking, tool routing, parsing, citation building, and an
evidence check --- work that must complete before citations and confidence can
be known.

\subsection*{Decision}

\texttt{POST /query/stream} computes the complete answer first, then emits it
as Server-Sent \texttt{content} events split on whitespace boundaries
(eight words per chunk by default), followed by one \texttt{metadata} event
carrying provider, latency, route, confidence, sources, retrieval scores,
citations, recommended actions, and the guardrail decision. Each event's
payload is JSON-encoded onto a single \texttt{data:} line.

The source states the reason for both choices: the split ``keeps every
separator, so concatenating the chunks reproduces the original text byte for
byte'', and JSON encoding is used ``so answer text containing newlines never
breaks the SSE framing''.

\subsection*{Consequences}

\begin{itemize}
  \item The endpoint behaves identically for every provider, including the
    local one, and is testable without network access.
  \item Clients can render progressively and then attach evidence, which the
    bundled frontend does.
  \item Time to first byte is not improved --- it equals the full answer
    latency --- so the endpoint delivers the streaming \emph{interface} without
    the streaming \emph{benefit}. The measured latency is reported in the
    metadata event, which makes this visible to clients.
  \item Newlines in answers cannot corrupt the event framing.
\end{itemize}

\subsection*{Alternatives}

True token streaming from providers that support it, with a fallback for those
that do not. That would make the endpoint's timing behaviour depend on
configuration and would require citations to be sent before the answer is
complete or omitted entirely.

\subsection*{Evidence}

\srcfile{src/feedback\_intelligence\_agent/api.py}
(\texttt{query\_stream}, \texttt{\_split\_for\_streaming}, \texttt{\_sse\_event},
\texttt{\_build\_stream\_metadata} and their docstrings);
\srcfile{src/feedback\_intelligence\_agent/schemas.py}
(\texttt{StreamMetadata}); \srcfile{frontend/src/api.ts};
\srcfile{tests/test\_api.py}.

\section{ADR-010: Nearest-rank percentiles for latency reporting}
\label{adr:010}

\subsection*{Status}

Accepted; in force at revision \texttt{\projectrevisionshort}.

\subsection*{Context}

The benchmark harness reports a 95th-percentile latency per pipeline phase.
Percentile definitions differ, and interpolating definitions can return a value
that was never observed.

\subsection*{Decision}

\texttt{benchmarking.percentile} implements the nearest-rank method:
$\rho = \lceil p n \rceil$ clamped to $[1, n]$, returning $y_{(\rho)}$ from the
sorted samples. The source gives the rationale directly: nearest-rank ``always
returns an actually-observed sample, which is easy to reason about for latency
SLOs and avoids fabricating values between measurements''.

\subsection*{Consequences}

\begin{itemize}
  \item Every reported percentile corresponds to a real measurement.
  \item The function is exact and has no floating-point interpolation, so it is
    trivially testable and reproducible.
  \item For small sample counts the estimate is coarse: with ten repetitions,
    $\operatorname{P}_{0.95}$ is the maximum. The report carries the sample
    count for every statistic, so the coarseness is visible rather than
    implied.
  \item Numbers are not directly comparable with tools that interpolate, such
    as NumPy's default \texttt{linear} method.
\end{itemize}

\subsection*{Alternatives}

Linear interpolation between order statistics (NumPy's default), or a
streaming estimator such as t-digest. The first fabricates unobserved values;
the second is disproportionate for a local benchmark harness.

\subsection*{Evidence}

\srcfile{src/feedback\_intelligence\_agent/benchmarking.py}
(\texttt{percentile} and its docstring; \texttt{summarize} reporting mean,
median, p95, min, max, and sample count);
\srcfile{tests/test\_benchmarking.py}.

